\section{Correspondence with the Manuscript}
\label{app:correspondence}

This table identifies the manuscript location to revise and the corresponding
statement in this reference. Numbers in the left column refer to Overleaf
\texttt{0450d48}; numbers in the right column refer to this PDF. The following
appendix links each reference statement to its Lean declarations. Thus a reader
can follow a manuscript definition to the mathematical replacement and then,
if needed, to the checked source.

\begin{longtable}{@{}>{\raggedright\arraybackslash}p{.32\textwidth}>{\raggedright\arraybackslash}p{.63\textwidth}@{}}
\toprule
Manuscript & Reference and relationship \\
\midrule
\endhead
Section 3, configuration and Apply/Merge rules &
Definition~\ref{def:execution} retains $C=\langle S,H,P,\visw\rangle$ and the
same transition updates, making rank, clock and fresh-fork premises explicit.
Definition~\ref{def:issuance} adds original issuance evidence. \\[5pt]
Lemma 3.1, LCA &
Definition~\ref{def:gca} fixes the intended \emph{greatest} ancestor;
Theorem~\ref{thm:gca-intersection} supplies the event-origin hypothesis behind
the intersection argument. Definition~\ref{def:virtual} is an extension for
multiple merge bases. \\[5pt]
Section 4, Notations &
Definition~\ref{def:events} keeps event replay and the operation policy.
Definition~\ref{def:invariant} adds a justified invariant to the commutation
domain; implementation equality remains concrete. \\[5pt]
Definitions 4.1--4.2, linearization relation and linear extension &
Definition~\ref{def:order} keeps the two clauses and local absorber;
the same invariant domain is used in both noncommutation tests. \\[5pt]
Definition 4.3, sequential specification &
Definition~\ref{def:history} retains the language generated by a labelled
transition system and allows an explicit event-to-label map. Specification
conflict is an additional definition on that language. \\[5pt]
Section 4.1, three replay laws &
Definition~\ref{def:guards} makes event guards explicit. Definitions
\ref{def:issuance} and \ref{def:scoped-laws} state the distinct certified
alternative; its scope must not be omitted. \\[5pt]
Definition 4.4, RA-linearizability &
Definition~\ref{def:ra} is a substantive extension: one witness must also
preserve specification-conflicting visibility. The preceding paragraph
states the literal manuscript criterion separately. \\[5pt]
Theorem 4.5; Definitions 4.6--4.7, canonical states &
Section~\ref{sec:semantics}'s canonical-state discussion and
Theorem~\ref{thm:scoped-replay} distinguish global linear-extension uniqueness
on supported, causally closed sets from scoped legal-replay uniqueness.
Existence is an explicit obligation. \\[5pt]
Section 5.1, WeakClosed, downset, Admissible, VC1--VC5 &
Definitions~\ref{def:representation}--\ref{def:replay-supply} state the
representation, metadata dependencies, equations and reconstruction supplies.
These are hypotheses, not consequences of query agreement. \\[5pt]
Lemma 5.1, Join &
Theorem~\ref{thm:join} proves the union representation by peeling.
Canonical join follows after relating representation to canonical replay. \\[5pt]
Theorem 5.2, canonical correctness &
Theorem~\ref{thm:execution} lifts a canonical join contract to certified
executions and recursive virtual bases. It states the internal replay order
separately from the public history order. \\[5pt]
Section 6, broken set; Theorem 6.1 and its simulation premise &
Example~\ref{ex:broken-set} distinguishes literal RA from the strengthened
criterion. Theorem~\ref{thm:sequential-lift} gives the literal lifting step;
Equation~\eqref{eq:commutation-compatibility} supplies a sufficient extra
condition for specification visibility. The simulation uses the paper's
specification-refinement function $\alpha$. \\[5pt]
Extension for guarded datatypes &
Definition~\ref{def:history-adequacy} and Theorem~\ref{thm:history-bridge}
replace unrestricted fold soundness by a certified history argument with
explicit concrete-state uniqueness. Examples~\ref{ex:queue}--\ref{ex:fugue} show its
scope and limitations. \\
\bottomrule
\end{longtable}

\paragraph{Presentation changes and mathematical changes.}
The section order, notation and names in this revision follow the manuscript.
The transition system and datatype implementations have not been changed by
this rewrite. The invariant commutation domain, original issuance evidence,
metadata reconstruction premises and specification-visibility clause are
mathematical differences from the manuscript, already present in the checked
baseline. They require explicit manuscript changes rather than a renaming.
The specification function $\alpha$ in Section 6 relates concrete states to
states of the independent sequential specification.
